\documentclass[11pt]{article}
\usepackage[letterpaper,margin=1in]{geometry}
\usepackage{amsmath,amssymb,amsthm}
\usepackage{graphicx}
\usepackage{float}
\usepackage[hidelinks]{hyperref}
\usepackage{microtype}

\newtheorem{theorem}{Theorem}
\newtheorem{lemma}{Lemma}
\newtheorem{corollary}{Corollary}
\newcommand{\C}{\mathbb C}
\newcommand{\Z}{\mathbb Z}

\title{A nontrivial weakly amenable Beurling algebra\\
on a noncommutative group\\
\large Full answer to the existence challenge in arXiv:1702.06605}
\author{Automated mathematics run \texttt{fa\_banach\_001}\\
Lane 5, agent \texttt{agent\_lane\_05}}
\date{9 August 2026}

\begin{document}
\maketitle

\begin{center}
\fbox{\parbox{0.91\textwidth}{\textbf{Status.}
Candidate full solution, likely valid, pending human review.  The example has
an unbounded polynomial weight and a genuinely noncommutative convolution
algebra.}}
\end{center}

\section{Source challenge}

Shepelska and Zhang write in the introduction of \cite{SZ}, printed page 3
(PDF page 3), that weak amenability of noncommutative Beurling algebras was
still unclear and that they had not seen a nontrivial weakly amenable example
which was noncommutative.  This restates the first open problem in
Zhang's earlier paper \cite{Zhang}.

The relevant sentence is:
\begin{quote}
``So far we have not even seen a non-trivial example of weakly amenable
Beurling algebra which is not commutative.''
\end{quote}
We give such an example explicitly.

\section{Definitions}

For a discrete group $G$ and a weight $\omega:G\to[1,\infty)$ satisfying
$\omega(xy)\leq\omega(x)\omega(y)$, the Beurling algebra
$\ell^1(G,\omega)$ is the convolution algebra with norm
\[
       \|f\|_\omega=\sum_{x\in G}|f(x)|\omega(x).
\]
A Banach algebra $A$ is \emph{weakly amenable} if every bounded derivation
$A\to A^*$ is inner.  It is \emph{amenable} if every bounded derivation from
$A$ into the dual of every Banach $A$-bimodule is inner.

We write $A\widehat\otimes B$ for the projective tensor product Banach
algebra.  The dual module conventions are
\[
 \langle c,a\cdot\Phi\rangle=\langle ca,\Phi\rangle,
 \qquad
 \langle c,\Phi\cdot a\rangle=\langle ac,\Phi\rangle.
\]

\section{Proof intuition}

The construction adds a finite noncommutative direction to a familiar
commutative weakly amenable Beurling algebra.  The main point is to prove that
this does not create a new derivation.

For a derivation on $A\widehat\otimes B$, amenability of $B$ first removes
its restriction to the $B$-factor by subtracting one inner derivation.  The
remaining restriction to $A$ takes values in the functionals which are
central for $B$.  Because $A$ is commutative, that space is a symmetric
$A$-module.  A short duality argument shows that unital commutative weak
amenability kills derivations into every such module.  Hence the original
derivation is inner.

Taking $A=\ell^1(\Z,(1+|n|)^\alpha)$ and $B=\ell^1(S_3)$ completes the
example: the first algebra is weakly amenable for $\alpha<1/2$, while the
second is amenable by its explicit finite-group diagonal.

\begin{figure}[H]
  \centering
  \includegraphics[width=\textwidth]{figures/open_problem_crop.png}
  \caption{The source context and existence challenge, rendered from
  \texttt{source\_paper.pdf}, PDF page 3.}
\end{figure}

\section{The tensor theorem}

We first record the elementary module lemma needed in the proof.

\begin{lemma}\label{lem:symmetric}
Let $A$ be a unital commutative weakly amenable Banach algebra.  If $X$ is a
unital symmetric Banach $A$-bimodule, then every bounded derivation
$\delta:A\to X$ is zero.
\end{lemma}

\begin{proof}
Fix $\lambda\in X^*$.  Define $T_\lambda:A\to A^*$ by
\begin{equation}\label{eq:slice}
  \langle T_\lambda(a),c\rangle
      =\langle\lambda,c\cdot\delta(a)\rangle
      \qquad(a,c\in A).
\end{equation}
Since $X$ is symmetric, for $a,b,c\in A$ we have
\begin{align*}
 \langle T_\lambda(ab),c\rangle
 &=\langle\lambda,c\cdot(a\cdot\delta(b)+\delta(a)\cdot b)\rangle\\
 &=\langle\lambda,(ca)\cdot\delta(b)+(cb)\cdot\delta(a)\rangle\\
 &=\langle a\cdot T_\lambda(b)+T_\lambda(a)\cdot b,c\rangle.
\end{align*}
Thus $T_\lambda$ is a bounded derivation $A\to A^*$.  Weak amenability makes
it inner; since $A$ is commutative, every inner derivation into $A^*$ is zero.
Putting $c=1_A$ in \eqref{eq:slice} gives
$\langle\lambda,\delta(a)\rangle=0$.  This holds for every $\lambda\in X^*$,
so $\delta(a)=0$ for every $a$.
\end{proof}

\begin{theorem}\label{thm:tensor}
Let $A$ be a unital commutative weakly amenable Banach algebra and let $B$ be
a unital amenable Banach algebra.  Then
$A\widehat\otimes B$ is weakly amenable.
\end{theorem}

\begin{proof}
Put $C=A\widehat\otimes B$ and let $D:C\to C^*$ be a bounded derivation.
Regard $A$ and $B$ as the commuting closed subalgebras
$A\otimes 1_B$ and $1_A\otimes B$ of $C$.

The restriction
\[
       b\longmapsto D(1_A\otimes b)
\]
is a derivation from $B$ into the dual $B$-bimodule $C^*$.  Amenability of
$B$ supplies $\Phi\in C^*$ such that this restriction agrees with the inner
derivation $c\mapsto c\cdot\Phi-\Phi\cdot c$.  Subtract that inner derivation
from $D$.  It is enough to prove that the resulting derivation $D_0$ is zero,
and by construction
\begin{equation}\label{eq:Bzero}
                D_0(1_A\otimes b)=0\qquad(b\in B).
\end{equation}

Let
\[
 X=\{\Psi\in C^*:(1_A\otimes b)\cdot\Psi
                  =\Psi\cdot(1_A\otimes b)\text{ for every }b\in B\}.
\]
This is a closed unital $A$-submodule of $C^*$.  It is symmetric as an
$A$-module because $A\otimes1_B$ is central in $C$: commutativity of $A$
gives
\[
        (a\otimes1_B)c=c(a\otimes1_B)\qquad(a\in A,\ c\in C).
\]

Define $\delta:A\to C^*$ by $\delta(a)=D_0(a\otimes1_B)$.  The two factors
commute, so applying $D_0$ to
$(a\otimes1_B)(1_A\otimes b)=(1_A\otimes b)(a\otimes1_B)$ and using
\eqref{eq:Bzero} yields
\[
  \delta(a)\cdot(1_A\otimes b)
     =(1_A\otimes b)\cdot\delta(a).
\]
Hence $\delta$ is a bounded derivation $A\to X$.  Lemma
\ref{lem:symmetric} gives $\delta=0$.

Finally, on every elementary tensor,
\[
 D_0(a\otimes b)
 =D_0\big((a\otimes1_B)(1_A\otimes b)\big)=0.
\]
Elementary tensors are dense in $C$, so $D_0=0$.  Therefore $D$ is inner,
as required.
\end{proof}

\section{Full explicit answer}

\begin{corollary}\label{cor:example}
For every $0<\alpha<1/2$, let
\[
   G=\Z\times S_3,
   \qquad
   \Omega(n,\sigma)=(1+|n|)^\alpha.
\]
Then $G$ is noncommutative, $\Omega$ is an unbounded symmetric weight, and
the Beurling algebra $\ell^1(G,\Omega)$ is weakly amenable and
noncommutative.  In particular it is a nontrivial weakly amenable Beurling
algebra of the type requested in \cite{SZ,Zhang}.
\end{corollary}

\begin{proof}
The inequality
\[
  1+|m+n|\leq(1+|m|)(1+|n|)
\]
shows that $\Omega$ is a weight.  It is symmetric and unbounded, so it is not
equivalent to the constant weight.  The group $G$ is noncommutative because
$S_3$ is noncommutative.

Set
\[
 A=\ell^1(\Z,\omega_\alpha),\qquad
 \omega_\alpha(n)=(1+|n|)^\alpha,
 \qquad B=\ell^1(S_3).
\]
Bade--Curtis--Dales proved that $A$ is weakly amenable exactly when
$0\leq\alpha<1/2$ \cite{BCD}; this is also recalled explicitly in the source
paper \cite[p.~3]{SZ}.  The finite group algebra $B$ is amenable.  Indeed,
\[
       M=\frac1{|S_3|}\sum_{\sigma\in S_3}
              \delta_\sigma\otimes\delta_{\sigma^{-1}}
       \in B\widehat\otimes B
\]
is a diagonal: $b\cdot M=M\cdot b$ for $b\in B$ and the multiplication map
sends $M$ to $\delta_e$.  Thus $B$ is even contractible.

The canonical convolution map gives the isometric Banach-algebra
identification
\[
  A\widehat\otimes B
     \cong \ell^1(\Z\times S_3,\Omega).
\]
Theorem \ref{thm:tensor} now gives weak amenability.  Finally, convolution is
noncommutative because the point masses supported on two noncommuting
elements of $S_3$ do not commute.
\end{proof}

The example is also a standard polynomial-weight example rather than merely
an arbitrary inflation: for any finite generating set of $\Z\times S_3$, its
word length differs from $|n|$ by a uniformly bounded amount in the finite
$S_3$ direction.  Thus $\Omega$ is equivalent to the corresponding usual
polynomial word-length weight.

\section{Verification and scope}

The proof has no computational or unproved dependency.  Its main checkable
points are:
\begin{enumerate}
\item amenability of $B$ is used with the dual $B$-module $C^*$, exactly as
      allowed by the definition;
\item after the $B$-restriction is killed, commutation of the two tensor
      factors forces the $A$-restriction into the $B$-central submodule;
\item centrality of $A\otimes1_B$ makes that module symmetric;
\item Lemma \ref{lem:symmetric} is proved directly from the defining
      $A\to A^*$ form of weak amenability;
\item $\Omega$ is unbounded and the algebra is genuinely noncommutative.
\end{enumerate}

The result fully answers the source's existence challenge.  It does not settle
the source's broader tensor-product question asking whether the projective
tensor product of \emph{arbitrary} weakly amenable Beurling algebras must be
weakly amenable.  Theorem \ref{thm:tensor} gives a rigorous positive class:
one factor may be any unital amenable algebra, provided the other is unital,
commutative, and weakly amenable.

\section{Bounded novelty check}

On 9 August 2026 we searched the run's registry, solution, attempt, and
proof-gap indexes; the local arXiv source corpus; exact phrases from the
source challenge; and web searches for noncommutative weakly amenable
Beurling algebras and weak amenability of projective tensor products.  No
independent verified source for Corollary \ref{cor:example} was found.

A later arXiv preprint of Mehdipour and Rejali \cite{MR}, Corollary 4.3,
states a much broader pullback conclusion which would imply noncommutative
examples.  Its proof depends on Proposition 4.1(i).  There, from injectivity
of a pullback map $\mathfrak S:Q(G_2)\to Q(G_1)$ and the inclusion
$\mathfrak S(I(G_2,\omega))\subseteq I(G_1,\overleftarrow\omega)$, the
argument asserts that an arbitrary non-inner quasi-additive function upstairs
has the form $\mathfrak S(p_2)$.  This is a surjectivity assertion not supplied
by the preceding hypotheses or proof.  Consequently that preprint does not,
without a repair, provide a verified duplicate.  The construction here uses
a different, complete amenable-factor argument.  Novelty confidence is
therefore moderate and remains subject to expert review.

\section{Human-review recommendation}

Review as a candidate full affirmative answer to the existence challenge.
The main verifier focus should be Theorem \ref{thm:tensor}, especially the
passage to the $B$-central symmetric module and the slicing argument in Lemma
\ref{lem:symmetric}.  A secondary literature review should check whether this
amenable-factor tensor lemma, or this exact $\Z\times S_3$ construction, has
appeared elsewhere under a different formulation.

\begin{thebibliography}{9}
\bibitem{BCD}
W.~G. Bade, P.~C. Curtis Jr., and H.~G. Dales,
\emph{Amenability and weak amenability for Beurling and Lipschitz algebras},
Proc. London Math. Soc. (3) \textbf{55} (1987), 359--377.

\bibitem{MR}
M.~J. Mehdipour and A. Rejali,
\emph{Weak amenability of weighted group algebras},
arXiv:2209.08346v1 (2022).

\bibitem{SZ}
V. Shepelska and Y. Zhang,
\emph{Non-weakly amenable Beurling algebras},
arXiv:1702.06605 (2017); Indiana Univ. Math. J. \textbf{67} (2018), 119--150.

\bibitem{Zhang}
Y. Zhang,
\emph{Weak amenability of commutative Beurling algebras},
Proc. Amer. Math. Soc. \textbf{142} (2014), 1649--1661;
arXiv:1207.5033.
\end{thebibliography}

\end{document}
